%----------
%	EDA - TARGET VARIABLES
%----------	


\cite{lameiras-2024} provides a detailed explanation of the methodology for the textual analysis of hate speech in the aforementioned corpus. Different categories have been defined specifically for this work based on previous research \cite{calderon-2020}. In this case, relevant categories for this thesis are "Análisis General" (AG), "Contenido Negativo" (CNEG) and "Insultos" (INS). 

%\begin{table}[H]
%    \centering
%    \begingroup
%    \small % Change to desired font size
%    \ttabbox[\FBwidth]{
%        \caption{Categories and Variables of the Analysis}
%        \label{tab:data_categories} % Etiqueta única para la tabla
%    }{
%        \begin{tabularx}{\textwidth}{
%            >{\hsize=1.2\hsize}X
%            >{\hsize=0.8\hsize}X
%            >{\hsize=1.2\hsize}X
%            >{\hsize=0.8\hsize}X
%        }
%            \hline
%            \rowcolor{gray!30} % Adding a gray color to the header row
%            \textbf{Category} & \textbf{Data Column} & \textbf{Variables} & \textbf{Source} \\
%            \hline
%            Categoría primaria de análisis de comentarios & Análisis General & Positivo, negativo o neutro & Calderón et al., (2020); Ada Lameiras (2019) \\
%            \hline
%            Categoría secundaria de análisis de comentarios negativos & Contenido Negativo & Estereotipos de género, sexualización/objetivación, violencia de género y Desprestigio/infravaloración que se dividió en: desprestigio hacia la víctima, desprestigio hacia la deportista autora de la publicación e infravaloración del acto o la acción & Ada Lameiras (2019) \\
%            \hline
%            Categoría terciaria de análisis de las diferentes formas de proferir insultos & Insultos & Genéricos, sexistas/misóginos, sarcasmo, imposición, objetivación sexual o deseo de dañar & Crosas-Remón y Medina-Bravo (2018); Piñeiro-Otero y Martínez-Rolán (2021); Cividanes-Álvarez y Martínez Rolán (2023) \\
%            \hline
%            \multicolumn{4}{l}{Source: \cite{lameiras-2024}} \\
%            \hline
%        \end{tabularx}
%    }
%    \endgroup
%\end{table}


In particular, the following definitions have been provided by Professor Alba Adá for the classes that were finally considered after the preprocessing detailed in Section~\ref{sec:data_preprocessing}.

\textbf{Análisis General}
\begin{itemize}
    \item Comentario Positivo: "Información o declaración de cualidades positivas y/o logros, que reconocen y valoran de forma favorable a una persona/empresa/acto/movimiento".\\
    (eng.) \textit{Information or statement of positive qualities and/or achievements, which favourably recognize and value a person/company/act/movement.}
    
    \item Comentario Negativo: "Información o declaración desfavorable y/o crítica que rechaza y expresa descontento hacia una persona/empresa/acto/movimiento".\\
    (eng.) \textit{Unfavorable and/or critical information or statement that rejects and expresses dissatisfaction towards a person/company/act/movement.}
\end{itemize}

\textbf{Contenido Negativo}
\begin{itemize}
    \item Desprestigiar Víctima: "Comentarios que se centran en dañar la reputación de la víctima (en este caso Jennifer Hermoso)".\\
    (eng.) \textit{Comments that focus on damaging the reputation of the victim (in this case Jennifer Hermoso).}
    
    \item Desprestigiar Acto: "Información que se basa en desacreditar lo sucedido, quitarle importancia al acto del beso no consentido".\\
    (eng.) \textit{Information based on discrediting what happened, downplaying the importance of the act of non-consensual kissing.}
    
    \item Desprestigiar Deportista Autora: "Comentarios que se centran en dañar la reputación de la persona que escribe el tuit".\\
    (eng.) \textit{Comments that focus on damaging the reputation of the person writing the tweet.}
    
    \item Insultos: "El usuario/a usa un lenguaje ofensivo, abusivo y vulgar dirigido sistemáticamente a una persona o grupos de personas".\\
    (eng.) \textit{The user uses offensive, abusive and vulgar language systematically directed at a person or groups of people.}
\end{itemize}

\textbf{Insultos}
\begin{itemize}
    \item Sexistas/Misóginos: "Lenguaje ofensivo y denigrante dirigido específicamente a las mujeres por el simple hecho de serlo o a su comportamiento sexual".\\
    (eng.) \textit{Offensive and demeaning language directed specifically at women simply because they are women or at their sexual behaviour.}
    
    \item Genéricos: "Lenguaje ofensivo que se utiliza para menospreciar, humillar y de manera ofensiva hacia una persona o grupos de personas de una manera amplia".\\
    (eng.) \textit{Offensive language that is used to belittle, humiliate and offend a person or groups of people in a broad manner.}
    
    \item Deseo de Dañar: "El usuario/a se centra en el deseo de provocar daño de manera explícita en otras personas".\\
    (eng.) \textit{The user is focused on the desire to explicitly cause harm to other people.}
\end{itemize}




